\begin{abstract}


Este proyecto estudia la estructura de una población de exoplanetas confirmados del catálogo PSCompPars del NASA Exoplanet Archive mediante Mapper, una técnica de Topological Data Analysis aplicada sobre datos procesados e imputados. La motivación central no es producir una clasificación definitiva de planetas, sino explorar si las relaciones entre sus propiedades físicas, orbitales y térmicas generan regiones estructuradas que puedan ser interpretadas astrofísicamente. Al mismo tiempo, el proyecto reconoce que los datos de exoplanetas no son una muestra neutral: están afectados por valores faltantes, por decisiones de imputación y por sesgos observacionales asociados a los métodos de descubrimiento.

La pregunta que guía el análisis es si Mapper/TDA puede revelar regiones significativas en el espacio de propiedades de los exoplanetas, y hasta qué punto esas regiones reflejan señal física u orbital, sesgo observacional, dependencia de imputación o sensibilidad metodológica. Para responderla, se reconciliaron 21 grafos Mapper existentes, correspondientes a siete espacios de variables y tres lentes, todos construidos con la cubierta fija \texttt{cubes10\_overlap0p35}. Los espacios comparados incluyen variables físicas, orbitales, térmicas y combinadas. La interpretación principal se concentró en los grafos con lente \texttt{pca2}, mientras que los lentes \texttt{domain} y \texttt{density} se consideraron como capas de sensibilidad.

El grafo que emergió como candidato principal fue \texttt{orbital\_pca2\_cubes10\_overlap0p35}. Este Mapper contiene 124 nodos, 196 aristas y $\beta_1=96$, con una imputación media nodal de 0.011. Esta combinación lo vuelve especialmente interesante: muestra una estructura topológica rica, pero con baja dependencia de valores imputados. En contraste, el Mapper térmico también presenta complejidad, pero su imputación media nodal de 0.588 obliga a interpretarlo con mayor cautela.

Para evitar una lectura ingenua de la topología observada, se incorporó una auditoría de sesgo observacional. Las variables \texttt{discoverymethod}, \texttt{disc\_year} y \texttt{disc\_facility} no se usaron para construir los grafos, sino como metadata externa para evaluar si las regiones detectadas estaban asociadas con la forma en que los planetas fueron descubiertos. En el grafo orbital principal, la pureza observada ponderada por método fue 0.842, el NMI observado entre asignación nodal y método de descubrimiento fue 0.254, y la prueba nula por permutación produjo $p=0.001$. Esto sugiere que la estructura orbital contiene información relevante, pero también está significativamente asociada con sesgos de descubrimiento.

En conjunto, los resultados muestran que Mapper puede identificar regiones candidatas útiles para inspección científica, especialmente en el espacio orbital. Sin embargo, esas regiones no deben interpretarse como la topología real de los exoplanetas ni como una taxonomía final. La conclusión principal es mixta: hay evidencia de estructura física u orbital plausible, pero también una contribución medible del sesgo observacional y de la incertidumbre introducida por la imputación. Por ello, la aportación del proyecto es una lectura crítica de la topología observada, separando regiones con posible señal astrofísica de regiones observacionales, mixtas o débiles.


\end{abstract}
